\bookchapter{Elements, isotopes and composite matter}
\label{sec:vii-elements}

\subsection{From a nuclear address to a periodic family}

The nuclear ledger distinguishes proton number, neutron number and the channels carried by their populations. A periodic construction asks a different question: how should the finite matter account organize families whose members share a proton number? The move from an isotope to an element family removes information for a particular purpose. It does not turn every isotope of the family into the same nuclear state.

The current source construction assigns one native family row to each integer $Z$ from $1$ through $126$. It audits the native coordinate before applying a familiar name or measured comparator. The order is important because an externally supplied label cannot count as a prediction of a construction that used the label to choose its own input. In this account the native row is constructed first; the label is a downstream field.\cite{viisama,viireview}

The finite periodic surface is useful as an organized model inventory. It places known-label families and source-native frontier families within one declared indexing scheme. Its value does not depend on calling every row an observed element. The scheme can be exact as an inventory while the physical status of an individual frontier row remains unestablished. The distinction is carried explicitly through this chapter.

Chemistry belongs to this volume because elements, isotopes and composites continue the matter construction. This placement does not imply that a complete molecular chemistry operator has been derived. This chapter supplies the family indexing, nuclear populations and source-support channels. The later chapters then develop the completed occupied atomic and molecular source constructions, including full pair fields, collective recoil and paid motion. The family inventory and the dynamical calculation retain their own definitions.

\subsection{Three row types and two maps}

An element-family row is denoted $E_Z$. An isotope row has the more informative address
\begin{equation}
I_{Z,N}=(Z,N,A_{\rm nuc}=Z+N,\ldots),
\end{equation}
where the remaining fields are defined by the source schema. The family projection is
\begin{equation}
\pi_Z(I_{Z,N})=E_Z.
\end{equation}
This map is many-to-one whenever the retained isotope roster contains more than one neutron count for the same proton count. The information it discards is not invalid information. It is simply not needed for the grouping question answered by the periodic surface.

A representative-isotope table requires another map,
\begin{equation}
\sigma_{\rm rep}(E_Z)=I_{Z,N_{\rm rep}(Z)}.
\end{equation}
The selector supplies one isotope for a one-row-per-element comparison. It does not invert the family projection on the entire isotope roster. Nor does it assert that the chosen isotope is the only allowed, stable or physically important member of the family. Those statements would require additional predicates.

The distinction can be followed directly in a readout. First construct $E_Z$. Then record the selector that supplies $N_{\rm rep}(Z)$. Next evaluate the nuclear or binding operator on the selected $(Z,N)$ pair. Finally compare the resulting field with a measured value if the row has an authorized comparator. If the selected isotope changes, the family coordinate may remain fixed while the readout changes. The selector is therefore part of the comparison specification.

This separation also explains why a periodic roster and an isotope vault can have different sizes without disagreement. The LC05 account retains 126 native family rows and an older scoped roster of 162 isotope rows. Neither number describes the cardinality of every possible isotope family. The first is the declared range of the periodic construction. The second is a particular retained source roster. Their coexistence is expected once the row types are stated.

\subsection{Native construction before label reveal}

For each permitted coordinate the source emits a row of the form
\begin{equation}
E_Z=(Z,{\rm native\_id}_Z,{\rm frontier}_Z,\ldots).
\end{equation}
The ellipsis is not permission to fill absent physical columns by analogy. It represents only fields already supplied by the source construction. In particular, a native identifier does not supply a shell configuration, a half-life, a decay channel or a molecular interaction law.

The periodic invariant audit checks that the coordinate lies in the declared range, that the number of rows is 126, and that the coordinates are unique. After that audit the label stage supplies 118 known labels for $Z\le118$ and eight native frontier identifiers for $119\le Z\le126$. These are descriptions of the source table's reveal policy. A successful label join validates the correspondence it actually checks; it does not by itself establish every physical property associated with the conventional name.

The chronological order prevents label-first backfilling. In a label-first procedure, a familiar name could select an isotope and the isotope could then be used to construct the row claimed to explain that name. That would make a visually convincing table without preserving the direction of inference. The retained procedure puts the invariant audit before reveal so that the native inventory can be inspected without relying on the display labels.

This does not make the full construction independent of all historical discovery. A method can be native-first in its executed replay while its design arose from earlier examination of familiar data. The stronger claim of prospective empirical prediction requires a corresponding frozen design and an eligible unseen comparison. The source order is a necessary piece of that discipline, not a substitute for the complete history of discovery and testing.

\subsection{The two surfaces of cardinality 126}

The retained matter capacity is $M=126$, and two major source surfaces share that count. Their partitions are different. The finite matter inventory is organized as
\begin{equation}
126=63\ \text{single}+63\ \text{bound},
\end{equation}
while the periodic surface is organized as
\begin{equation}
126=118\ \text{known labels}+8\ \text{frontier}.
\end{equation}
These decompositions make clear why equal cardinality cannot be used as an elementwise identity. The membership predicates differ, and the rows carry different kinds of addresses.

A row-position pairing between the two tables would certainly be possible as a display operation: first row with first row, second with second, and so on. Such a pairing is not a source-derived physical relation. To give it scientific meaning, the pairing would need to preserve a declared property or follow an independently specified map. The mere existence of a bijection between finite sets of equal size is insufficient.

The source account calls their common value a structural contact at the retained horizon. That statement can be retained without overextending it. The construction organizes both surfaces using a common finite capacity while allowing them to answer different questions. Keeping two schemas is not an admission of numerical inconsistency. It is what permits their respective information to remain available.

The same reasoning applies to other totals in the volume. The 321-row catalog, its 169 bound entries, the 80-row stable structural surface, and the 100-row computation grammar are not rival answers to one census question. They are different selected views. A periodic-family reader need not regard every carrier or rejected catalog candidate as an element, and an executable grammar need not contain a separate operation for every isotope represented in an input table.

\subsection{The reported periodic replay}

The LC05 source replay reports retention of all 126 native periodic rows, the 118 downstream labels, eight native frontier identities, the scoped 162-row isotope roster and the frontier seal inherited from CR071. Its decimal replay of the one-eighth/seven-eighths source split reports maximum error $10^{-96}$ against a $10^{-90}$ tolerance. These are recorded replay quantities, not a new execution undertaken for this manuscript.\cite{viireview}

The diagnostic alternatives include label backfill, row deletion, primitive mutation and frontier relabeling. Each challenges a specific part of the construction. Deleting a row challenges the declared inventory; changing a primitive challenges dependence on the locked structural inputs; replacing a frontier status with a known label challenges the distinction between native identity and external observation. Their purpose is to make clear what has been preserved by the replay.

A numerical tolerance in a decimal implementation is compatible with exact underlying rational identities. It records the implementation's arithmetic behavior. It does not enlarge the physical meaning of the identities. A very small split-reconstruction error is evidence about the split calculation, not evidence that a frontier nucleus has been synthesized or that its chemical behavior follows from the same tolerance.

The appropriate result is therefore a finite periodic and isotope record with reproducible labels and source custody. The completed replay establishes this inventory. The occupied electronic and molecular constructions developed later in this edition provide the subsequent dynamical recipients. 

\subsection{The balanced frontier addresses}

The eight source-native frontier families use balanced anchors
\begin{equation}
N=Z,\qquad A_{\rm nuc}=2Z.
\end{equation}
Their names are internal native identifiers of the SAM construction. They are retained here to make the source record identifiable, without presenting them as recognized observations.

\begin{table}[t]
\centering
\small
\begin{tabular}{lrrr}
\hline
Native identity & $Z$ & $N$ & $A_{\rm nuc}$\\
\hline
Harlium &119&119&238\\
Brockium &120&120&240\\
Uniquium &121&121&242\\
Liamium &122&122&244\\
Cooperium &123&123&246\\
Lindesium &124&124&248\\
Dorisium &125&125&250\\
Jerroldium &126&126&252\\
\hline
\end{tabular}
\caption{Source-native balanced frontier addresses. The table specifies model identities and arithmetic, not measured existence, stability or synthesis.}
\label{tab:vii-frontier}
\end{table}

At the terminal address,
\begin{equation}
Z=N=126=M.
\end{equation}
It follows that
\begin{equation}
A_{\rm nuc}=252=2M=14\Theta.
\end{equation}
The balanced nuclear ledger gives
\begin{equation}
u=2Z+N=378,\qquad d=Z+2N=378,
\end{equation}
and the neutral electron count is $e=Z=126$. With the source volume address $\hat V=54$,
\begin{equation}
u=d=7\hat V.
\end{equation}
Every equality in this chain follows from the displayed definitions. The chain does not require an empirical measurement of the frontier row in order to be an exact internal identity.

CR273 records a five-way terminus contact involving the matter horizon, the source doubly-magic predicate, carrier-mode inheritance, noble-cipher inheritance and component symmetry. These descriptions refer to the source predicates. A familiar phrase such as ``doubly magic'' or ``noble'' must not be allowed to import an unstated experimental nuclear or chemical conclusion. The arithmetic and model predicates remain available for inspection; observation, stability and chemical behavior remain separate statuses.

The balanced choice also has a clear domain. It supplies a particular anchor for each frontier family. It does not establish that the balanced isotope is the physically realized or longest-lived isotope, nor that all other neutron counts are excluded. The family and the anchor continue to be different objects even when the table displays one row per family.

\subsection{Historical readout and zero-refit extension}

The historical nuclear readout enters after the nuclear ledger and asymmetry columns have been constructed. CR274 reports a 55-row fitted benchmark. The base readout has RMS $3.9102$ MeV. Four fitted family gates reduce the reported RMS to $2.7160$ MeV, with 50 of 55 rows within 5 MeV. The fitted gates are parameters of that empirical construction. They do not become substrate primitives merely because the resulting expression uses structural coordinates.\cite{viireview}

CR277 then freezes that operator and extends it to one representative isotope per $Z$. It first reproduces the retained benchmark predictions to less than $10^{-3}$ MeV and then evaluates the extended table without refitting. The sequence is important: construction and fitting precede freezing; benchmark reproduction checks the implemented frozen operator; extension examines its behavior on a broader representative surface.

The reported observed-label subset contains 118 rows and has RMS $6.672$ MeV. Its light subset, $Z=1$ through $7$, has seven rows and RMS $19.907$ MeV. The middle subset, $Z=8$ through $82$, has 75 rows and RMS $4.145$ MeV. The heavy subset, $Z=83$ through $118$, has 36 rows and RMS $5.752$ MeV. The extended-only comparison has 78 rows and RMS $7.974$ MeV.

These figures should be read together. The full-table residual is not represented adequately by the more favorable middle subset, and the light-row miss should not disappear when a global value is quoted. The zero-refit extension provides information precisely because the operator was not adjusted to absorb every new deviation. It reveals a domain in which the frozen readout transfers less well.

Zero refit is a statement about the extension, not the origin of the coefficients. It does not convert the historical fitted operator into a wholly parameter-free law. Nor does it turn a representative-isotope evaluation into a test over all isotopes of all element families. The frozen readout, the representative selector, the comparator source and the reported residual must remain joined in the account.

This historical chain is also distinct from the later current B5T binding authority discussed in the preceding chapter. The tests have different representations, data splits and questions. Their RMS values should not be ranked as though they were interchangeable scores from one benchmark. The historical periodic readout remains evidence about that construction's extension; B5T remains a later representation and held-chain residual result with its own physical limits.

\subsection{What makes a composite more than a collection}

A collection of admitted constituents is the beginning of a composite description, not its completion. A composite also requires a connector pattern, a debit or credit operator, a closure gate and a source-emission rule. The source abstraction is
\begin{equation}
C=\bigl(\{\mathfrak r_i\},\Gamma,\mathcal D,
\mathcal G_{\rm close},\mathcal E_{\rm source}\bigr).
\end{equation}
The constituent set records which rows are present. The connector pattern $\Gamma$ records how they are related. The operator $\mathcal D$ specifies the relevant debit or credit. The gate decides whether that connected configuration closes. The emission rule states how an admitted composite returns support.

Adding catalog indices cannot recover these fields. Indices are identifiers, and changing the order of a table can change an index without changing the physical or grammatical object it identifies. Even if an index happened to equal a useful scalar, the sum of indices would still omit orientation, pair class, gate state and emitted support. The composite construction must preserve the relationships that make the collection one object.

This explains the different composite counts encountered earlier. CR252 retains 169 bound entries inside the 321-row catalog. LC06 retains 63 bound writes inside a 126-row matter inventory. A four-slot nucleus is yet another representation. The entries are selected and grouped under different predicates. A complete account can preserve all three without manufacturing a universal row-number correspondence.

The four-slot ledger supplies constituent population and channel information. It does not, by itself, specify every connector of a three-dimensional object. A graph-based reception construction may add exact incidence for a particular object, while an empirical binding operator supplies a separate energy readout. The fact that these descriptions can be joined at selected interfaces does not make each description independently sufficient for all the others.

\subsection{From periodic families to occupied and moving matter}

The periodic construction supplies family coordinates, labels, neutral electron counts and isotope grouping. The complete nuclear source adds conditional states and histories. The next step receives that source in an occupied electronic system.

The following chapters develop the completed scalar and polarization Hamiltonian, Pauli occupation, direct and exchange pairs, whole ion/free references and collective recoil. Their molecular continuation retains all 91 carbon and 153 oxygen moving reference channels, source-sector kinetic motion and paid virtual binding witnesses. These dynamical constructions receive the inventory and supply the additional state and interaction data. Reaction rates, resolved molecular spectra and environmental response retain their separate observables.\samref{pair_source}\samref{molecular_motion}

